\chapter{Contexte général et problématique}
\label{chap:contexte}

\section*{Introduction}
\addcontentsline{toc}{section}{Introduction}

Ce chapitre introduit le contexte scientifique et les motivations du présent
travail sur la détection des émotions multiples dans les textes. Il commence
par situer le cadre institutionnel du projet, avant de formuler la
problématique de recherche. Les objectifs et les questions de recherche qui
en découlent sont ensuite présentés, puis le chapitre se conclut par la
description du plan de travail adopté, qui structure la démarche
méthodologique suivie tout au long du mémoire.

\section{Cadre général du projet}
\label{sec:cadre-general}

Ce travail s'inscrit dans le cadre du Projet de Fin d'Études du Mastère de
Recherche en Informatique, spécialité Sciences et Technologies de
l'Information, réalisé à l'Institut Supérieur d'Informatique (ISI) de
l'Université de Tunis El Manar, sous l'encadrement du Dr.~Sahbi Bahroun. Le
sujet porte sur la conception d'une architecture neuronale pour la
détection des émotions multiples dans les textes, combinant attention
hiérarchique et connaissances sémantiques externes.

L'analyse automatique des émotions dans le texte s'est imposée comme un
sous-domaine central du traitement automatique des langues, avec des
applications directes en modération de contenu, en analyse de l'expérience
client, en santé mentale numérique et en systèmes conversationnels
affectifs. Contrairement à l'analyse de sentiment classique, qui réduit un
texte à une polarité (positif, négatif, neutre) \cite{pang2008,mohammad2020sentiment},
la détection d'émotions vise à identifier des catégories affectives fines
(joie, colère, tristesse, surprise, etc.), souvent multiples et simultanées
au sein d'un même énoncé.

\section{Problématique}
\label{sec:problematique}

La plupart des travaux en classification des émotions textuelles reposent
sur l'hypothèse simplificatrice qu'un texte exprime une émotion dominante
unique. Cette hypothèse ne reflète pas la réalité des
états affectifs humains, qui se manifestent fréquemment sous forme de
mélanges d'émotions coexistant dans un même énoncé, par exemple
l'expression simultanée de la surprise et de la joie, ou de la déception et
de la colère. Les jeux de données récents à annotation multi-étiquettes,
tels que GoEmotions \cite{demszky2020}, montrent qu'environ 17\,\% des
textes portent au moins deux émotions simultanées, ce qui invalide
l'approche mono-étiquette pour un grand nombre de cas réels.

Trois limites structurelles freinent les approches existantes de
détection multi-émotions, détaillées et documentées au chapitre~2 :

\begin{itemize}
\item \textbf{L'indépendance supposée des étiquettes.} Les architectures de
  classification multi-label standards traitent chaque émotion comme une
  décision binaire indépendante, ignorant les corrélations et
  incompatibilités entre émotions (la joie et la tristesse sont rarement
  co-actives, tandis que la surprise accompagne souvent la peur ou la
  joie).
\item \textbf{La perte de la structure du texte.} L'encodage global d'un
  texte en un unique vecteur de représentation efface les signaux
  localisés à l'échelle de la phrase, alors que dans les textes
  multi-phrases, différentes émotions sont souvent portées par des
  segments distincts du texte.
\item \textbf{Le déséquilibre sévère des classes.} La distribution des
  émotions dans les corpus réels suit une loi de puissance : quelques
  émotions fréquentes (joie, neutralité) dominent très largement les
  émotions rares (deuil, soulagement, nervosité), ce qui pénalise
  l'apprentissage des classes minoritaires, précisément celles pour
  lesquelles la détection est la plus critique en pratique.
\end{itemize}

Ces limites motivent la question de recherche centrale de ce mémoire :
\textit{dans quelle mesure une architecture combinant encodage
hiérarchique, attention croisée entre émotions et enrichissement par
connaissances sémantiques externes peut-elle améliorer la détection des
émotions rares et multiples dans le texte, par rapport à un modèle
pré-entraîné affiné de façon standard ?}

\section{Objectifs de recherche}
\label{sec:objectifs}

Pour répondre à cette question, ce travail poursuit les objectifs suivants :

\begin{itemize}
\item \textbf{Concevoir} une architecture modulaire, HAKE-MER, structurée
  en quatre composantes cumulables (encodage hiérarchique, attention
  croisée multi-émotions, enrichissement sémantique, pondération
  dynamique), permettant d'isoler la contribution de chaque mécanisme ;
\item \textbf{Évaluer} cette architecture de façon rigoureuse sur le jeu de
  données GoEmotions, par un protocole d'ablation cumulative à budget
  d'entraînement strictement constant, sur deux backbones pré-entraînés
  (DistilBERT et RoBERTa), afin de garantir que les gains observés sont
  attribuables aux composantes architecturales et non à des facteurs
  confondus ;
\item \textbf{Comparer} deux ressources de connaissances externes de nature
  différente (lexique catégoriel NRC contre ressource conceptuelle
  SenticNet) afin de déterminer laquelle s'intègre le mieux à
  l'architecture proposée ;
\item \textbf{Analyser} les résultats au niveau de la classe, au-delà des
  métriques globales, pour évaluer l'apport réel de chaque composante sur
  les émotions rares, enjeu applicatif central du sujet.
\end{itemize}

\section{Plan de travail}
\label{sec:plan-travail}

La démarche suivie dans ce mémoire s'inspire d'un processus itératif de
type CRISP-DM (Cross-Industry Standard Process for Data Mining), adapté au
contexte d'un projet de recherche en apprentissage profond. Le
Tableau~\ref{tab:plan-travail} résume les phases méthodologiques adoptées.

\begin{table}[H]
\centering
\caption{Phases méthodologiques du plan de travail}
\label{tab:plan-travail}
\begin{tabular}{|p{4cm}|p{9cm}|}
\hline
\textbf{Phase} & \textbf{Contenu} \\
\hline
Analyse du problème et revue de littérature &
Formulation de la problématique, revue critique des approches de
détection multi-émotions, d'attention hiérarchique et d'intégration de
connaissances externes (chapitre~2), identification du positionnement de
HAKE-MER. \\
\hline
Compréhension et préparation des données &
Exploitation des statistiques GoEmotions (chapitre~2), choix du
prétraitement et de la longueur de séquence. \\
\hline
Conception et modélisation &
Formulation du problème, conception de l'architecture HAKE-MER et de ses
quatre modules, définition des hypothèses de recherche H1 à H4
(chapitre~3). \\
\hline
Expérimentation et évaluation &
Implémentation, protocole d'ablation à budget constant (Step~0 sur
DistilBERT et RoBERTa ; ladder M1 à M4 sur DistilBERT), comparaison
NRC/SenticNet, analyse par classe (chapitre~4). \\
\hline
Documentation et valorisation &
Rédaction du mémoire, synthèse des enseignements méthodologiques et des
recommandations pratiques (chapitre~5), structuration de l'état de
l'art en survey autonome en vue d'une valorisation après soutenance. \\
\hline
\end{tabular}
\end{table}

\section*{Conclusion}
\addcontentsline{toc}{section}{Conclusion}

Ce premier chapitre a présenté le cadre général du projet, formulé la
problématique de la détection des émotions multiples dans les textes et
défini les objectifs de recherche qui en découlent, ainsi que le plan de
travail du mémoire. Pour rendre la question centrale opérationnelle, il
reste à fixer la tâche et le corpus de référence, à identifier une
baseline crédible et à examiner comment la littérature répond aux
limitations L1 à L3 (étiquettes, structure du texte, classes rares) avant
de positionner HAKE-MER ; c'est l'objet du chapitre~2.
